\section*{Hoofdstuk \ref{ch:b3:cytoskeleton-motility} --- Dynamiek van het cytoskelet en celbeweging}

\begin{solution}{exo:b3:cytoskeleton-motility:1}
Actine: \qty{7}{nm}, bolvormig actine, polair (gekarteld/gepunt), ATP;
cortex, uitstulping, samentrekking met \omterm{def:b3:cytoskeleton-motility:motors}{myosine}. \omterm{def:b3:cytoskeleton-motility:filaments}{Microtubulus}:
\qty{25}{nm}, $\alpha\beta$-tubulinedimeer, polair (plus/min), GTP;
transport over lange afstand, spoelfiguur, trilharen. \omterm{def:b3:cytoskeleton-motility:filaments}{Intermediair
filament}: \qty{10}{nm}, coiled-coil-eiwitten (keratine, vimentine,
lamine), apolair, geen nucleotide; mechanische sterkte van cel en kern.
\end{solution}

\begin{solution}{exo:b3:cytoskeleton-motility:2}
De monomeerconcentratie waarbij een uiteinde noch groeit noch krimpt,
$k_{\text{off}}/k_{\text{on}}$. De twee uiteinden verschillen omdat de
subeenheid haar ATP na de inbouw hydrolyseert, zodat de uiteinden
verschillende chemische soorten tonen (ATP-actine dat aan het
gekartelde uiteinde aankomt, ADP-actine dat het gepunte uiteinde
verlaat) met verschillende affiniteiten. Bij gelijke \omterm{thm:b3:cytoskeleton-motility:treadmilling}{kritische
concentraties} zou er één evenwicht zijn: geen treadmilling, geen
stroom, en geen manier om een gerichte omzet te betalen.
\end{solution}

\begin{solution}{exo:b3:cytoskeleton-motility:3}
Naar de rand over \omterm{def:b3:cytoskeleton-motility:filaments}{microtubuli}: \omterm{def:b3:cytoskeleton-motility:motors}{kinesine}, naar het plusuiteinde. Naar
het centrum: cytoplasmatisch \omterm{def:b3:cytoskeleton-motility:motors}{dyneïne}, naar het minuiteinde.
Spanningsvezel: \omterm{def:b3:cytoskeleton-motility:motors}{myosine} II, naar het gekartelde uiteinde van actine.
\omterm{def:b3:cytoskeleton-motility:cilia}{Trilhaar}: \omterm{def:b3:cytoskeleton-motility:cilia}{axonemaal dyneïne}, naar het minuiteinde van de naburige
dubbel (de basis).
\end{solution}

\begin{solution}{exo:b3:cytoskeleton-motility:4}
Uitstulping (Rac voor lamellipodia, Cdc42 voor filopodia), hechting
(\omterm{def:b3:cytoskeleton-motility:migration}{integrines}; stroomafwaarts van Rac en Rho), samentrekking (Rho, via
\omterm{def:b3:cytoskeleton-motility:motors}{myosine} II), intrekking van de achterkant (het loslaten van de
adhesies, aangedreven door de samentrekking).
\end{solution}

\begin{solution}{exo:b3:cytoskeleton-motility:5}
$C_{c}^{+} = 1/5 = \qty{0.2}{\micro M}$, $C_{c}^{-} = 0.8/1 =
\qty{0.8}{\micro M}$; $C_{\text{ss}} = (1 + 0.8)/(5 + 1) = \qty{0.3}{\micro
M}$; $J = 5\times 0.3 - 1 = 0.5$ subeenheden per seconde, \qty{1.35}{nm/s}.
\end{solution}

\begin{solution}{exo:b3:cytoskeleton-motility:6}
$800/8 = 100$ ATP per seconde. De meter kost $1/(8\times 10^{-7}) =
1.25\times 10^{6}$ s, ongeveer twee weken, en $1.25\times 10^{8}$ stappen
en ATP. Diffusie zou \num{16000} jaar kosten: een verschil van een
miljard keer duizend.
\end{solution}

\begin{solution}{exo:b3:cytoskeleton-motility:7}
$F_{\max} = 8.3\times 10^{-20}/3.6\times 10^{-8} = \qty{2.3}{pN}$.
Dezelfde energie over een langere stap geeft minder kracht --- een
langere hefboom. \omterm{def:b3:cytoskeleton-motility:motors}{Kinesine} (korte stap, \qty{6}{pN}) past bij zware
lasten; \omterm{def:b3:cytoskeleton-motility:motors}{myosine} V (lange stap) legt per ATP meer afstand af en past bij
snel transport van lichte vracht.
\end{solution}

\begin{solution}{exo:b3:cytoskeleton-motility:8}
(a) Taxol bevriest de \omterm{def:b3:cytoskeleton-motility:filaments}{microtubuli}: de spoelfiguur kan niet zoeken,
vangen en corrigeren, en de mitose loopt vast bij het controlepunt; de
kruipende cel blijft uitstulpen maar verliest haar polariteit op
termijn. (b) Colchicine haalt de \omterm{def:b3:cytoskeleton-motility:filaments}{microtubuli} weg: geen spoelfiguur,
mitotische stilstand; het kruipen gaat op actine door maar zonder
volgehouden richting. (c) Cytochalasine dekt de gekartelde uiteinden
af: geen uitstulping, geen kruipen; de mitose verloopt wel, maar de
cytokinese faalt, wat tweekernige cellen geeft. (d) Remming van Rac:
geen lamellipodia, dus geen kruipen; de deling blijft grotendeels
onaangetast.
\end{solution}

\begin{solution}{exo:b3:cytoskeleton-motility:9}
Verblijfstijd van de orde van de halveringstijd, \qty{20}{s}
(tijdconstante $20/\ln 2 \approx \qty{29}{s}$); onbeweeglijke fractie
\qty{10}{\%}. De polymerisatie aan de rand moet zowel de voortgang als
de achterwaartse stroom leveren: $1 + 1.5 = \qty{2.5}{\micro m/min}$.
\end{solution}

\begin{solution}{exo:b3:cytoskeleton-motility:10}
Een groeiend uiteinde houdt een kap van GTP-tubuline; daarachter maakt
de hydrolyse gespannen GDP-tubuline. Het verlies van de kap --- een
toevallige gebeurtenis waarvan de kans stijgt naarmate de groei
vertraagt --- laat de spanning los en het uiteinde krimpt snel tot er
een nieuwe kap ontstaat. Net boven de \omterm{thm:b3:cytoskeleton-motility:treadmilling}{kritische concentratie} is elke
\omterm{def:b3:cytoskeleton-motility:filaments}{microtubulus} op elk moment ofwel voorzien van een kap en groeiend,
ofwel zonder kap en instortend, zodat de lengtes uiteenlopen in een
lange populatie en een verdwijnende in plaats van zich rond een
gemiddelde te scharen. De cel wint er een snelle verkenning van haar
volume mee en het vermogen om selectief te stabiliseren --- een
plusuiteinde dat een kinetochoor of een plek in de cortex bereikt wordt
gevangen en gehouden, de overige worden binnen minuten hergebruikt:
zoeken en vangen.
\end{solution}

\begin{solution}{exo:b3:cytoskeleton-motility:11}
Op deze schaal is de traagheid verwaarloosbaar en zijn de
vloeistofvergelijkingen omkeerbaar in de tijd: een beweging die
zichzelf terugvolgt (een roeispaan die uitgaat en terugkomt) brengt het
lichaam terug waar het begon, hoe snel elke slag ook is. Een draaiende
helix volgt zichzelf nooit terug --- haar beweging is chiraal en
doorlopend --- en levert dus netto stuwkracht. Bij \qty{100}{Hz} met
duizend protonen per omwenteling is dat $10^{5}$ protonen per seconde
per motor.
\end{solution}

\begin{solution}{exo:b3:cytoskeleton-motility:12}
De trilharen van de luchtwegen kunnen niet slaan, slijm en bacteriën
worden niet afgevoerd, en infecties keren terug. De zweepstaarten van
zaadcellen, op hetzelfde \omterm{def:b3:cytoskeleton-motility:cilia}{axoneem} gebouwd, zijn onbeweeglijk:
onvruchtbaarheid. In het embryo drijven beweeglijke trilharen van de
knoop een stroming naar links die de links-rechtsas vastlegt; zonder
die stroming wordt de kant willekeurig gekozen, zodat bij de helft van
de patiënten de organen zijn omgekeerd.
\end{solution}

\begin{solution}{pb:b3:cytoskeleton-motility:1}
\textbf{1.} $\qty{300}{\micro m^3} = 3\times 10^{-13}$ L; $\times
2\times 10^{-4}$ mol/L $= 6\times 10^{-17}$ mol $= 3.6\times 10^{7}$
moleculen; $1.8\times 10^{7}$ in filamenten.
\textbf{2.} $1.8\times 10^{7}\times \qty{2.7}{nm} = \qty{4.9}{cm}$
filament; bij \qty{0.25}{\micro m} per stuk ongeveer $2\times 10^{5}$
filamenten.
\textbf{3.} Afdekeiwit blokkeert de meeste gekartelde uiteinden, en
profiline en thymosine-$\beta$4 binden de monomeren, zodat het werkelijk
vrije ATP-actine bij de \omterm{thm:b3:cytoskeleton-motility:treadmilling}{kritische concentratie} ligt; de groei blijft
beperkt tot de uiteinden die de cel vrijmaakt.
\textbf{4.} $C_{\text{ss}} = 2.2/12.9 = \qty{0.17}{\micro M}$; $J = 11.6
\times 0.17 - 1.4 = 0.6$ subeenheden per seconde; \qty{1.6}{nm/s} $=
\qty{0.1}{\micro m/min}$.
\textbf{5.} $1/0.1 = 10$-voudig.
\textbf{6.} \qty{1}{\micro m/min} $= \qty{16.7}{nm/s} = 6.2$ subeenheden
per seconde per filament; $\times 2\times 10^{5} = 1.2\times 10^{6}$ ATP
per seconde, ongeveer $0.1\,\%$ van de omzet van de cel.
\textbf{7.} $1/(8\times 10^{-7}) = 1.25\times 10^{6}$ s $\approx 14.5$
dagen; $1.25\times 10^{8}$ stappen en ATP.
\textbf{8.} $L^{2}/2D$: $(1)^{2}/(2\times 10^{-12}) = 5\times 10^{11}$ s
$\approx \num{16000}$ jaar voor een meter; $(10^{-5})^{2}/(2\times
10^{-12}) = 50$ s voor \qty{10}{\micro m}. Diffusie voldoet binnen een
cellichaam; alles wat langer is vergt motoren.
\textbf{9.} $5\times 10^{4}/6.02\times 10^{23} = 8.3\times 10^{-20}$ J
$= 20\,k_{B}T$.
\textbf{10.} $8.3\times 10^{-20}/8\times 10^{-9} = \qty{10}{pN}$;
rendement $6/10 \approx 60\,\%$.
\textbf{11.} $F = 6\pi\times 0.01\times 10^{-7}\times 8\times 10^{-7} =
1.5\times 10^{-14}$ N $= \qty{0.015}{pN}$, vierhonderd keer onder de
\omterm{prop:b3:cytoskeleton-motility:physics}{blokkeerkracht}: één motor is ruim voldoende tegen de viskeuze weerstand;
de echte lasten zijn obstakels en verankeringen.
\textbf{12.} $10^{4}\times 100 = 10^{6}$ ATP per seconde, $0.1\,\%$ van
de begroting van het neuron: het transport is goedkoop. Maar motoren
stoppen zodra het ATP daalt, zodat het falen van de mitochondriën de
synaps berooft van alles wat in het cellichaam wordt gemaakt --- het
falende axonale transport dat bij neurodegeneratieve ziekten wordt
gezien.
\textbf{13.} \qty{10}{\micro m/min} $= \qty{167}{nm/s}$; $167/2.7 = 62$
subeenheden per seconde per filament.
\textbf{14.} $200\times 20 = 4000$ filamenten; $4000\times 62 = 2.5\times
10^{5}$ subeenheden per seconde.
\textbf{15.} $2.5\times 10^{5}$ ATP per seconde, $0.025\,\%$ van de
begroting.
\textbf{16.} $\qty{3}{\micro m}/(\qty{10}{\micro m/min}) = \qty{0.3}{min}
= \qty{18}{s}$; $2.5\times 10^{5}\times 18 = 4.5\times 10^{6}$
subeenheden in het netwerk.
\textbf{17.} $1\,\text{pN}\times 2.7\,\text{nm} = \qty{2.7}{pN.nm} =
0.66\,k_{B}T$: een schommeling van die energie treedt op met kans
ongeveer $e^{-0.66} \approx 0.5$ --- spleten van één subeenheid gaan
voortdurend open, en de ratel is volstrekt aannemelijk.
\textbf{18.} De polymerisatie zelf duwt, waarbij de bacterie het
``membraan'' aan de voorkant van haar eigen netwerk is; met de
Arp2/3-machinerie van de cel haalt zij de snelheid van een \omterm{def:b3:cytoskeleton-motility:migration}{lamellipodium}
en meer --- tot \qty{0.5}{\micro m/s}, \qty{30}{\micro m/min}.
\textbf{19.} Bij $C = K_{d}$ is $\theta = 1/2$ en $\mathrm{d}\theta/
\mathrm{d}C = 1/(4C)$: een verschil van \qty{2}{\%} in $C$ geeft $\Delta\theta
= 0.005$; met \num{25000} receptoren per kant zijn er vooraan $125$ meer
bezet.
\textbf{20.} Elke kant heeft ongeveer \num{12500} bezette receptoren, die
schommelen met $\sqrt{\num{12500}} \approx 110$, zodat het verschil van
twee kanten met ongeveer $160$ schommelt: de $125$ verdwijnt op elk
moment in de ruis. Receptoren binden ongeveer eens per seconde opnieuw;
middelen over $N$ seconden krimpt de ruis met $\sqrt{N}$ --- tien
seconden brengen haar op $50$ en de gradiënt wordt gelezen.
\textbf{21.} Rac (met Cdc42) vooraan, Rho achteraan. Met overal actief
Rac stulpt de cel rondom uit, spreidt zij zich en komt zij nergens.
\textbf{22.} $30/3 = \qty{10}{\micro m/min}$: zoals gegeven.
\textbf{23.} De uitstulping is vervangen door blebbing dat door druk
wordt gedreven; hechting en samentrekking hangen nog steeds van actine
en \omterm{def:b3:cytoskeleton-motility:motors}{myosine} af, zodat actine voor de rest van de cyclus onmisbaar
blijft.
\textbf{24.} Het front wordt elke \qty{18}{s} opnieuw gebouwd: een
verschuiving in de activiteit van Rac stuurt de polymerisatie binnen één
levensduur van het netwerk om, het nieuwe front gaat voorop, en de
trage stappen achteraan volgen.
\textbf{25.} Treadmilling in vitro \qty{1.6}{nm/s}; ongeveer $14.5$
dagen om het axon af te leggen; zo'n $2.5\times 10^{5}$ ATP per seconde
voor de uitstulping.
\end{solution}
